% !Mode:: "TeX:UTF-8"

\chapter{绪论}

\section{研究背景与意义}

随着信息技术的飞速发展，各个领域（尤其是医疗诊断、金融风控等）积累了海量的高维数据。这些数据往往具有随机性、不完备性和不一致性，如何从中挖掘出有价值的知识并用于决策，是人工智能领域的重要挑战。在医疗领域，如乳腺癌等疾病的诊断数据，其特征（如肿块厚度、细胞大小等）不仅维度高、存在冗余，而且很多特征在数值上具有连续性，在临床意义上具有模糊性，这给疾病的精准分类和可解释性诊断带来了困难。

粗糙集理论（Rough Set Theory）作为一种处理模糊和不确定性数据的数学工具，由 Pawlak 于 1982 年提出\upcite{pawlak1982rough}。该理论建立在分类机制的基础上，无需依赖先验概率分布或人为设定隶属度，而是直接利用数据本身的内在结构进行推理。在六种决策一致性假设下，粗糙集通过“分类（粒化）”和“逼近”来研究具有随机性、不完备性和不一致性的信息系统，从而有效解决属性约简、异常检测、规则提取以及分类决策等问题。将粗糙集理论应用于乳腺癌医疗数据集，不仅能够剔除冗余特征、降低模型复杂度，还能提取出符合医学逻辑的规则，具有重要的理论意义与临床应用价值。

\section{国内外研究现状}

粗糙集理论建立在对知识的相对形式化基础之上。所谓知识，即描述了对象在属性下的值的集合。粗糙集理论中的信息系统通常形式化为四元组 \(\langle U, A, V, f \rangle\)，其中 \(U\) 为对象集（论域），\(A\) 为属性集，\(V\) 为值域，\(f\) 为映射函数。当属性集 \(A\) 中包含决策属性时，其余属性被称为条件属性（记为 \(C\)），决策属性记为 \(D\)，此时构成决策系统。

\subsection{经典粗糙集理论}
Pawlak 提出的经典粗糙集为后续研究指明了方向。其核心思想是：先选定属性集 \(B\) 并构造等价关系 \(R_B\) 对信息系统中的对象进行分类；然后，在近似空间中通过构造上下近似来逼近信息系统中的任意概念（即 \(U\) 的子集）。在决策系统中，算法会根据决策属性 \(D\) 将论域划分为多个决策类 \(D_i\)，分别逼近这些类并合并出正域和负域。将正负域从论域 \(U\) 中剔除后，剩余部分即为边界域（边界样本）。

有研究表明，直接在边界样本上学习得到的分类边界与在原始样本上学习的结果基本一致，且能大幅降低学习时间\upcite{hu2008neighborhood}。这种基于边界样本的检测方法，不仅提升了计算效率，也与异常点（离群点）检测有着深刻的内在联系。

\subsection{粗糙集理论的拓展}
经典粗糙集理论在早期为了简化模型，主要针对定性属性进行设计。然而现实世界中的数据极为复杂，从运算性质上可分为定性、定序、定比和定距属性。\textbf{经典粗糙集使用的等价关系无法有效利用数值属性所具备的良好运算性质}。此外，经典粗糙集的鲁棒性相对较差，对数据中的噪声较为敏感。

为解决上述问题，学术界提出了多种粗糙集泛化模型：
\begin{enumerate}
	\item \textbf{邻域粗糙集（Neighborhood Rough Sets）}：为了充分刻画和利用数值属性，邻域粗糙集引入距离度量代替等价关系。当规定邻域半径 \(\delta=0\) 时，邻域粗糙集便退化为经典粗糙集\upcite{hu2008neighborhood}。邻域半径的大小不仅反映了人们对数值数据噪声的容忍度，也反映了同类对象在数值属性上的真实差异尺度。
	\item \textbf{模糊粗糙集（Fuzzy Rough Sets）}：为了提升模型的鲁棒性并更好地处理模糊概念，学者们引入了模糊算子（如三角范数、三角余范数）和模糊等价关系，构建了模糊粗糙集模型\upcite{dubois1990rough}。它能够更加细腻地刻画数据的不确定性，而无需依赖脆弱的硬划分。
\end{enumerate}

\section{现有研究存在的问题}

尽管粗糙集理论已取得长足发展，但在实际应用中仍面临以下挑战：
\begin{enumerate}
	\item \textbf{离散化带来的信息损失}：经典粗糙集仅能处理离散化数据。在处理乳腺癌连续数值特征（如 WBC 数据集）时，必须进行离散化。然而，等宽或等频的离散化往往会破坏数据原有的连续运算性质，造成原始信息损失，从而影响后续约简和规则的准确性。
	\item \textbf{邻域半径的选取困境}：在邻域粗糙集中，虽然避免了离散化，但邻域半径 \(\delta\) 的合理设定极为困难。理论上，只要邻域半径足够小，信息系统总是分类一致的；但这会妨碍获得期望的知识粒子簇。当模型无法在合理大小的邻域半径上完成任务时，往往需要反复试错来去除冗余属性并补充条件属性，这增加了算法的不稳定性。
	\item \textbf{算法效率与鲁棒性的平衡}：如何在保证模型鲁棒性（如利用模糊粗糙集的模糊算子）的同时，设计出高效的属性约简与规则提取算法，仍是当前研究的重点。
\end{enumerate}

\section{本文主要研究内容}

针对上述问题，本文以医疗数据集（乳腺癌数据集 WBC）为背景，围绕粗糙集理论在属性约简与规则提取中的应用展开研究。主要工作包括：
\begin{enumerate}
	\item \textbf{基于经典粗糙集的预处理与约简算法设计}：针对 WBC 数据集的连续属性，对比不同离散化策略的影响。基于依赖度函数，设计并实现核属性求解（Core）与启发式属性约简（Reduction）算法，剔除冗余特征。
	\item \textbf{决策规则提取与应用}：在属性约简的基础上，定义规则的支持度与置信度，提取高可解释性的 IF-THEN 决策规则，并探讨其医学意义。
	\item \textbf{经典与模糊粗糙集的理论对比研究}：结合模糊粗糙集的理论优势，探讨其在处理连续数值属性及提升鲁棒性方面的潜能，为后续直接处理连续数据提供理论支撑。
\end{enumerate}

\section{论文组织结构}

本论文共分为五章，具体安排如下：
\begin{itemize}
	\item \textbf{第一章 绪论}：阐述研究背景与意义，总结粗糙集及其扩展模型的国内外研究现状，指出当前存在的问题，并介绍本文的主要研究内容。
	\item \textbf{第二章 粗糙集理论基础}：详细推导经典 Pawlak 粗糙集的数学定义（等价关系、上下近似、依赖度等），并引入模糊粗糙集的基础概念（三角范数、模糊等价关系）。
	\item \textbf{第三章 属性约简与规则提取算法设计}：重点介绍数据离散化策略、基于依赖度的核与约简算法逻辑，以及规则提取的置信度框架。
	\item \textbf{第四章 实验设计与结果分析}：在 WBC 数据集上验证算法有效性，展示约简前后的依赖度变化，分析提取的核心诊断规则，并对比经典与模糊理论的特性。
	\item \textbf{第五章 总结与展望}：总结全文工作，结合粒度计算与异常检测方向，对未来的研究进行展望。
\end{itemize}